\uuid{foPv}
\titre{Intégration/Quotient/DL usuels}
\auteur{hanani ; mzari}
\organisation{exo7}
\licence{CC BY-NC-SA 4.0}
\datecreate{2020-01-20}
\niveau{L1}
\module{Analyse}
\chapitre{Développement limité}
\sousChapitre{Calculs}
\difficulte{3}
\theme{Opérations sur les DL, Intégration, Quotient, DL usuels}
\exo7qcmid{lille-2.developpements-limites.025}

\contenu{
\texte{
Soit $a\in \Rr$ et $f$ une fonction. On dit que $f(x)$ est un $o\left((x-a)^n\right)$ au voisinage de $a$ et on écrit $f(x)=o((x-a)^n)$ si $\displaystyle \lim _{x\to a}\frac{f(x)}{(x-a)^n}=0$. On note par $DL_n(a)f(x)$ le développement limité de $f$ à l'ordre $n$ en $a$. On note aussi $DL_n(a^+)f(x)$ (resp. $DL_n(a^-)f(x)$) le développement limité de $f$ à l'ordre $n$ à droite en $a$ (resp. à gauche en $a$).
}

\begin{question}
\qclassification{125.02}
\qtags{motcle=Intégration/Quotient/DL usuels}
Soit $\displaystyle f(x)=\int _0^x\frac{\ln (1+t)}{\sqrt{1+t}}\mathrm{d}t$ et $\displaystyle g(x)=\frac{\ln (1+x)}{\sqrt{1+x}}$. Alors
\begin{answers}
\good{$\displaystyle DL_2(0)g(x)=x-x^2+o(x^2)$.}
\bad{$f'(x)$ ne possède pas de $DL(0)$.}
\bad{$f(x)$ ne possède pas de $DL(0)$.}
\good{$\displaystyle DL_3(0)f(x)=\frac{x^2}{2}-\frac{x^3}{3}+o(x^3)$.}
\end{answers}
\begin{explanations}
On a : $\displaystyle f'(x)=g(x)$ et $DL_2(0)g(x)=x-x^2+o(x^2)$. Par intégration, et puisque $f(0)=0$, on obtient :
$$DL_3(0)f(x)=\frac{x^2}{2}-\frac{x^3}{3}+o(x^3).$$
\end{explanations}
\end{question}
}
